\begin{figure}[b]
\centering
\resizebox{\columnwidth}{!}{%
\begin{tikzpicture}[
    >=Stealth,
    line width=1pt,
    font=\large,
    oma/.style={
        draw=black,
        rounded corners=3pt,
        minimum width=5cm,
        minimum height=2.7cm,
        align=center,
        line width=1.2pt
    }
]

% Central OMA box
\node[oma] (oma) {OMA};

% Input text
\node[left=3.0cm of oma] (input) {
    $\mathbf{X}$
};

% Output labels
\node[right=3.0cm of oma, yshift=1.1cm] (freq) {
    $f_k$
};

\node[right=3.0cm of oma] (damp) {
    $\zeta_k$
};

\node[right=3.0cm of oma, yshift=-1.1cm] (mode) {
    $\bm{\phi}_k$
};

% Input arrow
\draw[->]
    (input.east) -- (oma.west);

% Output arrows
\draw[->]
    (oma.east |- freq) -- (freq.west);

\draw[->]
    (oma.east) -- (damp.west);

\draw[->]
    (oma.east |- mode) -- (mode.west);

\end{tikzpicture}%
}
\caption{
    Schematic diagram of an OMA algorithm. 
    Input measurements organized in a $d\times T$ (measurement channels $\times$ time 
    points) matrix are fed into the OMA algorithm, which produces an estimate of the 
    natural frequencies $f_k$, damping ratios $\zeta_k$ and mode shapes $\bm{\phi}_k$.
    }\label{fig:oma-schematic}
\end{figure}